\section{How wrong is the face-value fit? A sensitivity bound}
\label{sec:sensitivity}

Let $s(r; \theta) = \partial_\theta \log \int_{c(r)} g(y \mid \theta)\, d\nu(y)$ be
the face-value score (the gradient of a summand of \eqref{eq:facevalue}) and
$\Info(\theta) = \E_\theta[\, s\, s^\top]$ the face-value information at $\delta =
0$. The analyst fits the face-value model regardless of $\delta$, so $\hat\theta$
solves $\sum_i s(R_i; \hat\theta) = 0$ while the reports are drawn from the
$\delta$-tilted law of \cref{def:tilt}. This is a misspecified-MLE problem
\citep{huber1967behavior,white1982maximum}: $\hat\theta \to_p \theta_\delta$, the
solution of $\E_\delta[s(R; \theta_\delta)] = 0$, and $\theta_0 = \theta^\star$ is
the truth.

\begin{theorem}[Graceful degradation under C2 violation]
\label{thm:sensitivity}
Assume C1, C3, and $\mathrm{C2}(\delta)$ with tilt $h \in \mathcal{H}_\delta$, and
the standard regularity that makes the face-value model a smooth, identifiable
model at $\delta = 0$ with nonsingular $\Info(\theta^\star)$. Then the face-value
pseudo-true parameter satisfies, as $\delta \to 0$,
\begin{equation}
  \theta_\delta - \theta^\star
  \;=\; \delta\, \Info(\theta^\star)^{-1}\,
        \Cov_{\theta^\star}\!\big(s(R; \theta^\star),\, h(R, Y)\big)
  \;+\; O(\delta^2),
  \label{eq:bias}
\end{equation}
where the covariance is taken over the joint law of $(Y, R)$ at $\delta = 0$. Because
$s$ is a function of $R$ alone, the tower property gives $\Cov(s, h(R,Y)) = \Cov(s,
\bar h)$ with $\bar h(R) = \E[h(R,Y) \mid R]$ the candidate-set-conditional tilt, so
the tilt enters the bias only through its candidate-set-conditional mean (the form
the appendix uses). In particular the asymptotic bias is bounded by
\begin{equation}
  \|\theta_\delta - \theta^\star\|
  \;\le\; \delta\, B(\theta^\star) + O(\delta^2),
  \qquad
  B(\theta^\star) = \big\|\Info(\theta^\star)^{-1}\big\|
     \sup_{\|h\|_\infty \le 1}
       \big\|\Cov_{\theta^\star}(s, h)\big\|,
  \label{eq:Bdelta}
\end{equation}
with $B(\theta^\star) < \infty$ and $B$ independent of $\delta$. The supremum is
attained at the worst-case tilt, the $L^\infty$ sign extreme point of the
candidate-set-conditional score component.
\end{theorem}

\begin{proof}[Sketch]
Write $\Psi(\theta, \delta) = \E_\delta[s(R; \theta)]$, so $\theta_\delta$ solves
$\Psi(\theta_\delta, \delta) = 0$ and $\Psi(\theta^\star, 0) = 0$ because the
face-value model is correctly specified at $\delta = 0$. The implicit function
theorem gives $d\theta_\delta/d\delta|_0 = -[\partial_\theta \Psi]^{-1}
\partial_\delta \Psi$ at $(\theta^\star, 0)$. The $\theta$-derivative is the
negative information, $\partial_\theta \Psi(\theta^\star, 0) = -\Info(\theta^\star)$
(information equality at the well-specified point). For the $\delta$-derivative,
differentiate the tilted report law $p_\delta(r \mid \theta) \propto \int_{c(r)}
g(y \mid \theta)\, \kappa_0(r) e^{\delta h(r,y)}\, d\nu(y)$: $\partial_\delta \log
p_\delta|_0$ is the within-candidate-set centered tilt, so $\partial_\delta
\Psi(\theta^\star, 0) = \E_{\theta^\star}[\, s(R; \theta^\star)\, \bar h(R)\,] =
\Cov_{\theta^\star}(s, \bar h)$, where $\bar h(r) = \E[h(R,Y) \mid R = r]$ is the
candidate-set-conditional tilt and we use $\E[s] = 0$ at $\delta = 0$. Combining
yields \eqref{eq:bias}; the bound \eqref{eq:Bdelta} is the supremum of
$\|\Cov(s, h)\|$ over $\|h\|_\infty \le 1$, attained at the $L^\infty$ sign extreme
point $\bar h = \operatorname{sign}$ of the candidate-set-conditional score
component, with the Cauchy--Schwarz step giving only a looser closed-form envelope
(deferred to \cref{app:sensitivity}). Full regularity (uniform integrability of the
remainder, a dominating envelope for the family $\{p_\delta\}$) is standard
$M$-estimation and is also deferred to \cref{app:sensitivity}.
\end{proof}

\begin{corollary}[Partial identification]
\label{cor:partial}
Under \cref{thm:sensitivity}, the identified set for $\theta^\star$ given the
coarsened data and a tilt budget $\delta$ is, to first order, the zonotope
\begin{equation}
  \mathcal{S}_\delta
  \;=\; \Big\{\, \hat\theta - \delta\, \Info^{-1}\Cov(s, h)
        \;:\; \|h\|_\infty \le 1 \,\Big\} + o(\delta),
  \label{eq:identset}
\end{equation}
the image of the $L^\infty$ tilt ball under the linear map $h \mapsto -\delta\,
\Info^{-1}\Cov(s, h)$, a centrally symmetric set centered at the face-value estimate
and contained in the ball of radius $\delta B(\theta^\star)$. (It is a zonotope, not
an ellipsoid: an ellipsoid would arise from an $L^2$ tilt constraint, whereas
\cref{def:tilt} bounds $h$ in $L^\infty$.) The set $\mathcal{S}_\delta$ contracts to
the point $\{\theta^\star\}$ as $\delta \to 0$, recovering the point identification
of \citet{towell2026synthesis} under C2.
\end{corollary}

\begin{remark}[What the marginal-fit check cannot see]
\citet{towell2026synthesis} shows that under C2 the fitted mean of $\T$ equals its
empirical mean, so a marginal-fit check is structurally blind to bias in the latent
parameter. \Cref{thm:sensitivity} quantifies the blind spot: the bias direction
$\Info^{-1}\Cov(s, h)$ is exactly the component of parameter error that leaves the
fitted marginal of $\T$ unchanged, so a larger marginal fit is never reassurance
against it. In a regular exponential family with natural statistic $\T$, when the
tilt is linear in $\T$ (that is $h = a^\top \T$) the leading term is exact along the
curved directions (the range of $\partial \eta / \partial \theta$) at every $\delta$,
because the report law stays an exponential family under the tilt, only with a
$\delta$-shifted natural parameter $\eta(\theta) + \delta a$; so \eqref{eq:bias} is
not merely a small-$\delta$ approximation there (see \cref{app:sensitivity}). The
location-family regime of \citet{towell2026synthesis} carries the usual
$O_p(n^{-1/2})$ remainder on top.
\end{remark}

The bound is only useful if $\delta$ can be pinned or bypassed. It cannot be
estimated from the coarsened data, so the next section supplies the external
device, the singleton, that both bounds $\delta$ where it matters and restores point
identification, and asks how much of it is needed.
